\section{Measurement Methodology}
\label{sec:methodology}

Every number in this report was produced by one command
(\texttt{scripts/run\_tournament.sh}, \cref{sec:reproducibility}) on one
host, in one continuous run, after the equivalence gate passed. This section
documents the host, the runtimes, the metrics, the workload, and the timing
boundaries, so that each figure can be traced to a script, a configuration
and a revision.

\subsection{Host and runtimes}
\label{sec:host}

\begin{table}[htbp]
    \centering
    \caption{Host and runtime versions of the run that produced the committed
    snapshot. The full record, including the filesystem and the MongoDB
    build, is \texttt{benchmarks\_results/environment.txt}.}
    \label{tab:host}
    \begin{tabular}{@{}ll@{}}
        \toprule
        Component & Value \\
        \midrule
        CPU & Intel Core i7-10750H @ 2.60 GHz, 12 logical cores \\
        RAM & 8 GB \\
        OS / kernel & WSL2 (Ubuntu userland), Linux 6.18.33.2-microsoft-standard-WSL2 \\
        Python & 3.14.4 (+ \texttt{psutil} for RSS/CPU) \\
        Go & 1.27.2 (\texttt{linux/amd64}) \\
        Rust & 1.99.0 (\texttt{rustc}/\texttt{cargo}, stable) \\
        Java (harness) & OpenJDK 17.0.20.1, executed via \texttt{mvn exec:java} \\
        Java (JMH) & OpenJDK 21.0.12.1, shaded jar \\
        Maven & 3.9.12 \\
        MongoDB & 8.0.13, \texttt{mongod} on \texttt{localhost:27017} \\
        \bottomrule
    \end{tabular}
\end{table}

All implementations run single-threaded and synchronous
(\cref{sec:architecture}); no benchmark phase uses worker pools, async I/O or
parallel iterators. The only other process involved is \texttt{mongod}, a
separate server whose CPU and memory are \emph{not} part of the client
measurements (\cref{sec:index-memory}).

\subsection{Metric definitions}
\label{tab:metrics}

\begin{table}[htbp]
    \centering
    \caption{What each column of the result CSVs means. Identical in all four
    languages; the harness enforces the schema
    (\texttt{scripts/verify\_format.py}).}
    \label{tab:metrics}
    \begin{tabularx}{\textwidth}{@{}lX@{}}
        \toprule
        Column & Definition \\
        \midrule
        \texttt{elapsed\_seconds} & wall-clock time of one operation, from a
        monotonic clock (\texttt{time.perf\_counter}, \texttt{time.Now},
        \texttt{Instant::now}, \texttt{System.nanoTime}); per-repetition
        builds report the arithmetic mean. \\
        \texttt{stdev\_seconds} & sample standard deviation of the per-operation
        elapsed times across repetitions (query latency: across all samples). \\
        \texttt{repetitions} / \texttt{iterations} & number of timed samples
        behind the row. \\
        \texttt{memory\_mb} & \textbf{RSS delta}: resident set size of the
        benchmark process \emph{after} minus \emph{before} the operation, in
        MB, read from \texttt{/proc/self/status} (Go/Rust/Java) or
        \texttt{psutil} (Python). It measures growth caused by the operation,
        not the total footprint of the process. \\
        \texttt{peak\_rss\_mb} & largest resident set size of the benchmark
        process observed while the operation ran, sampled every
        \fixed{50} ms by a background thread plus one final sample; across
        repetitions the maximum is reported. \\
        \texttt{cpu\_percent} & process CPU time (user + system) divided by
        wall-clock time, $\times 100$; above 100 means more than one core was
        saturated. \\
        \bottomrule
    \end{tabularx}
\end{table}

Both memory series describe the \emph{benchmark client process only}. For the
MongoDB structure, the server's own memory and disk are reported separately
(\cref{sec:index-memory}); they are never folded into the client's RSS.

\subsection{Corpus and workload}
\label{sec:workload}

\begin{table}[htbp]
    \centering
    \caption{Workload actually executed by every language.}
    \label{tab:workload}
    \begin{tabularx}{\textwidth}{@{}lX@{}}
        \toprule
        Parameter & Value \\
        \midrule
        Corpus & \fixed{100} Project Gutenberg books, tracked in
        \texttt{data\_source/} (bodies $\approx$\fixed{59} MB, headers
        $\approx$0.4 MB); book IDs listed in
        \texttt{benchmarks\_results/equivalence\_report.json}. \\
        Query workload & \fixed{30} active queries from the shared file
        \texttt{src/utils/benchmarks/queries.txt} (one query per line,
        terms separated by a single space; comments start with \texttt{\#}).
        Semantics: a document matches when it contains every term; an unknown
        term yields an empty result. \\
        Query samples & \fixed{5} repetitions $\times$ \fixed{30} queries =
        \fixed{150} timed samples per structure per language. \\
        Datalake lookup & \fixed{100} books located once per layout (header
        and body read). \\
        Download probe & \fixed{10} books downloaded sequentially over one
        connection; no concurrency anywhere in the pipeline. \\
        Build batches & \fixed{25, 50, 250, 500, 1000, 5000, 10000} books;
        5 repetitions up to 1000, 3 for 5000, 2 for 10000. \\
        Update batches & base indexes of \fixed{25, 50, 250, 500, 1000}
        books, each extended with \fixed{5} new books, \fixed{5}
        repetitions per base size (\cref{sec:index-updates}). \\
        Recovery & interrupt after \fixed{50} of \fixed{100} books, then
        resume; report processed, duplicated and lost IDs plus detection and
        processing time. \\
        \bottomrule
    \end{tabularx}
\end{table}

The scalability builds reuse the corpus by repetition: a ``10{,}000-book
batch'' is the \fixed{100}-book corpus indexed repeatedly until the batch
size is reached, so batch size measures index size, not corpus diversity.
The update experiment is the exception: it uses fresh synthetic book IDs
(\texttt{base\_size}$\times$100 + offset), so added books are genuinely new
terms and never collide with the base.

\subsection{Warm-up}
\label{sec:warmup}

Only the Java harness warms up, because Java is the only runtime with a JIT
compiler: before timing, each Java structure performs one unmeasured build of
the smallest batch and discards its artifacts, so the measured builds run on
already-compiled code. Python, Go and Rust are measured exactly as they
start: no warm-up, no discarded round. JMH numbers
(\cref{sec:jmh}) use JMH's own warm-up iterations and are reported
separately; they are never mixed with the cross-language figures.

\subsection{Cache, disk and network state}
\label{sec:environment-state}

\begin{itemize}
    \item \textbf{Artifact directories.} Every timed phase starts from a
    freshly cleaned artifact tree: the previous run's datalake, index and
    result files are deleted before the phase begins, so no phase reads a
    cache written by another.
    \item \textbf{OS page cache.} The tournament does not drop the system
    page cache (that requires root and would perturb the host). Files read
    twice within one phase may therefore be served from RAM. This affects all
    languages equally within a phase and is listed as a limitation
    (\cref{sec:limitations}).
    \item \textbf{Disk.} All artifacts live on the same filesystem as the
    repository (WSL2 / NTFS-backed). \texttt{disk\_usage.csv} reports file
    counts, directory counts and bytes for each layout and each index
    structure after storing the same corpus.
    \item \textbf{Network.} Only the download probe touches the network: ten
    sequential GETs, one connection at a time, measured with the same clock
    as everything else. Every other phase is fully offline and reads the
    tracked corpus from disk.
\end{itemize}

\subsection{Timing boundaries}
\label{sec:timing-boundaries}

\begin{itemize}
    \item \textbf{Write throughput} times the storage call itself: split
    header+body written to the layout (the download is timed separately), one
    book per iteration, \fixed{100} iterations per layout.
    \item \textbf{Download throughput} starts before the HTTP request and
    stops after the response body is fully read and written to disk.
    \item \textbf{Lookup} starts before locating the book's files under the
    layout and stops after both files are read into memory.
    \item \textbf{Index build} starts before the first token is produced and
    stops after the structure is fully written (JSON flushed, per-term files
    closed, or MongoDB insert acknowledged).
    \item \textbf{Query} starts before the postings lookup of the first term
    and stops after the intersection of all terms is computed.
    \item \textbf{Incremental update} starts before the first of the five new
    books is tokenized and stops after the structure is durably updated.
\end{itemize}

\subsection{Repetition and uncertainty}
\label{sec:uncertainty}

Each timed operation is repeated as specified in \cref{tab:workload}; plots
show the mean and, where a deviation band appears, it is the sample standard
deviation across those repetitions (\texttt{stdev\_seconds}). Readers can
assess uncertainty from the raw \texttt{repetitions} and \texttt{stdev}
columns committed next to every chart. No result is discarded, no outlier is
trimmed, and no phase is run twice to pick the better run: the committed
snapshot is a single continuous tournament.

\subsection{Coverage: what was compared with what}
\label{sec:coverage}

\begin{table}[htbp]
    \centering
    \caption{Language $\times$ structure coverage of the timed workload
    (D=datalake layouts, J=JSON index, F=folder index, M=MongoDB index,
    S=SQLite metadata).}
    \label{tab:coverage}
    \begin{tabular}{@{}lcccc@{}}
        \toprule
        Language & Datalake & JSON & Folder & MongoDB \\
        \midrule
        Python & D, S & J, F, M & & \\
        Go & D, S & J, F, M & & \\
        Rust & D, S & J, F, M & & \\
        Java & D, S & J, F, M & & \\
        \bottomrule
    \end{tabular}
\end{table}

Every language runs the full matrix: three datalake layouts plus the three
index structures, with identical parameters. MongoDB is skipped consistently
and identically by all four ports when the server is unreachable, so no
figure silently compares a MongoDB-less run with a MongoDB run.

\subsection{Revisions and provenance}
\label{sec:provenance}

Each chart in this report maps to (i) the script that produced it, (ii) the
raw CSVs in \texttt{benchmarks\_results/}, (iii) the configuration constants
listed in \cref{tab:workload}, and (iv) the Git revision recorded in
\texttt{benchmarks\_results/environment.txt} at run time. The mapping is
tabulated in \cref{sec:reproducibility}.
